The main confirmation contains 192 task--seed--protocol runs, representing
832 controller trajectories and 6,656 structural opportunities. The original
research record reports reconstruction checks of all completed histories for
actions, parameters, resource counts, and terminal risks. The archived research preparation separately replayed a representative
complete trajectory, as detailed in the historical reproducibility record. Figures~\ref{fig:learning} and
\ref{fig:decisions} display the full opportunity sequence, including holds.
Table~\ref{tab:main-results} reports every main proposal rule, rather than
selecting the best-performing rule per task.

\begin{table}[tb]
\centering\small
\begin{tabular}{llrrrr}
\toprule
Task & Proposal & $n_A=4096$ & $n_A=16384$ & Width & Grow $\ge2$ \\
\midrule
Null & Residual & 0.2532 & 0.2532 & 2.00 & 0.00 \\
 & Shuffled & 0.2532 & 0.2532 & 2.00 & 0.00 \\
 & Random & 0.2532 & 0.2532 & 2.00 & 0.00 \\
 & Covariance & 0.2532 & 0.2532 & 2.00 & 0.00 \\
 & Ssd & 0.2532 & 0.2532 & 2.00 & 0.00 \\
\addlinespace
Linear & Residual & 0.1720 & 0.1720 & 2.00 & 0.00 \\
 & Shuffled & 0.1720 & 0.1720 & 2.00 & 0.00 \\
 & Random & 0.1720 & 0.1720 & 2.00 & 0.00 \\
 & Covariance & 0.1720 & 0.1720 & 2.00 & 0.00 \\
 & Ssd & 0.1720 & 0.1720 & 2.00 & 0.00 \\
\addlinespace
Additive & Residual & 0.1472 & 0.1395 & 3.25 & 0.56 \\
 & Shuffled & 0.1472 & 0.1427 & 2.62 & 0.19 \\
 & Random & 0.1472 & 0.1377 & 3.50 & 0.62 \\
 & Covariance & 0.1472 & 0.1409 & 3.00 & 0.44 \\
 & Ssd & 0.1472 & 0.1381 & 3.38 & 0.62 \\
\addlinespace
Interaction & Residual & 0.1754 & 0.1507 & 5.62 & 1.00 \\
 & Shuffled & 0.1749 & 0.1531 & 5.25 & 1.00 \\
 & Random & 0.1739 & 0.1475 & 5.88 & 1.00 \\
 & Covariance & 0.1759 & 0.1498 & 5.69 & 1.00 \\
 & Ssd & 0.2113 & 0.2071 & 2.69 & 0.00 \\
\addlinespace
\bottomrule
\end{tabular}
\caption{Mean final Brier risk over 16 new seeds per task. Width and fraction
with at least two accepted growth events refer to the 16,384-label audit.
Both columns fit on the same 2,048 examples and use eight fresh audit blocks.
Uncertainty for the primary paired comparisons is in
Table~\ref{tab:paired-results}; complete marginal intervals are in the CSV.}
\label{tab:main-results}
\end{table}

\begin{figure}[tb]
\centering\includegraphics[width=\linewidth]{learning_trajectories.pdf}
\caption{Measured ordinary-start trajectories at 16,384 audit labels per
opportunity. Curves are mean retained risk and width; shaded bands are
pointwise 95\% Student-$t$ intervals across 16 independent seeds. Risk
integrates conditional label noise over 16,384 independent evaluation inputs.
All methods coincide on null and linear tasks. The same predictor persists
through a hold, while trial branches still consume computation.}
\label{fig:learning}
\end{figure}

\paragraph{Repeated growth and genuine no-growth.}
At the larger audit budget, the residual method accepts 20 additions across
16 additive-task runs and 58 across 16 interaction-task runs. Nine of 16
additive runs and all 16 interaction runs accept at least two additions.
Mean final widths are 3.25 and 5.625; final risks are
$0.13948\pm0.00280$ and $0.15067\pm0.00291$, where uncertainty denotes a
95\% interval half-width. Growth and holding occur within the same histories;
these are not restarts at independently constructed checkpoints. On both
null and linear tasks every primary finite-bound trajectory retains width
two and its warmup parameters. This prevents further overfitting but does
not recover Bayes risk automatically: null risk is 0.25325, above the known
0.25 Bayes risk because warmup itself was ungated.

At 4,096 audit labels, no residual additive run grows and each interaction
run grows exactly once, to width three. The interaction risk is 0.17538,
compared with 0.15067 at the larger audit budget. Fitting data and proposal
rules are fixed across this comparison, so the difference concerns the
evidence protocol and its induced history, not extra fitting samples.
It nevertheless costs four times as many audit labels. It is not a free
improvement in training efficiency.

\begin{figure}[tb]
\centering\includegraphics[width=\linewidth]{decision_histories.pdf}
\caption{Every primary residual trajectory at the larger audit budget.
Green means accepted growth; amber means accepted continuation; light gray
means the incumbent parameters are held. Rows are the 16 prespecified seeds,
columns the eight decision opportunities. An external opportunity schedule
exists, but actions and accepted final capacity are data-dependent.}
\label{fig:decisions}
\end{figure}

\paragraph{Proposal information is not a general advantage.}
Residual minus shuffled risk on the larger-audit additive task is
$-0.003268\pm0.002547$. This is a favorable comparison at equal candidate
optimization workload, with the caveat of descriptive multiple comparisons.
It does not survive the stronger interpretation of beating random proposal
addition: residual minus random is $+0.001813\pm0.002433$ on additive data
and $+0.003168\pm0.002128$ on interaction data. The latter descriptive paired interval excludes zero in random's favor while spending less on
candidate optimization. At the smaller audit budget, interaction residual
minus random is also adverse, $+0.001528\pm0.000854$.
The covariance rule is not reliably separated from the residual rule.
SSD is much worse on this interaction task, but its failure cannot establish
superiority over all growth principles: it perturbs a different restricted
action family and nearly fails to grow. These observations favor retaining
proposal alternatives instead of calling the selected residual proxy a
validated information principle.

\begin{table}[tb]
\centering\scriptsize
\begin{tabular}{lrrrr}
\toprule
Comparator & Null & Linear & Additive & Interaction \\
\midrule
shuffled & $+0.0000\pm0.0000$ & $+0.0000\pm0.0000$ & $-0.0033\pm0.0025$ & $-0.0024\pm0.0042$ \\
random & $+0.0000\pm0.0000$ & $+0.0000\pm0.0000$ & $+0.0018\pm0.0024$ & $+0.0032\pm0.0021$ \\
covariance & $+0.0000\pm0.0000$ & $+0.0000\pm0.0000$ & $-0.0014\pm0.0019$ & $+0.0009\pm0.0023$ \\
ssd & $+0.0000\pm0.0000$ & $+0.0000\pm0.0000$ & $+0.0014\pm0.0032$ & $-0.0565\pm0.0042$ \\
fixed matched & $-0.0022\pm0.0009$ & $-0.0007\pm0.0004$ & $-0.0008\pm0.0016$ & $-0.0005\pm0.0028$ \\
fixed matched tail & $-0.0007\pm0.0008$ & $-0.0004\pm0.0004$ & $-0.0008\pm0.0010$ & $-0.0023\pm0.0022$ \\
small & $-0.0016\pm0.0006$ & $-0.0001\pm0.0003$ & $-0.0066\pm0.0034$ & $-0.0595\pm0.0031$ \\
large & $-0.0229\pm0.0013$ & $-0.0176\pm0.0012$ & $-0.0120\pm0.0027$ & $-0.0067\pm0.0031$ \\
large pool & $+0.0030\pm0.0008$ & $+0.0014\pm0.0004$ & $+0.0080\pm0.0029$ & $+0.0104\pm0.0025$ \\
\bottomrule
\end{tabular}
\caption{Residual minus comparator Brier risk at the larger audit budget:
paired mean $\pm$ 95\% Student-$t$ half-width over 16 runs. Negative favors
residual growth. ``Fixed matched tail'' compares both predictors after 600
additional updates; other rows use final retained checkpoints. ``Large pool''
is width ten trained offline from the full label pool at matched training/probe
update proxy, with different minibatch optimization. Intervals are descriptive,
not multiplicity-adjusted discoveries.}
\label{tab:paired-results}
\end{table}

\begin{figure}[tb]
\centering\includegraphics[width=\linewidth]{attribution.pdf}
\caption{Selected paired comparisons from Table~\ref{tab:paired-results}.
Bars are 95\% intervals across seeds. The full-label-pool comparator changes
data allocation and optimizer schedule. The two tail curves continue training
without the gate and have no risk-monotonicity guarantee.}
\label{fig:attribution}
\end{figure}

\paragraph{Capacity and optimization controls limit the attribution.}
Compared with ungated width-two training, residual growth improves mean
larger-audit risk by 0.00655 on additive data and 0.05952 on interaction data.
Thus the implemented repeated process does realize useful new capacity on
these tasks. At retrospectively matched final capacity and training/probe
work, however, paired differences are
$-0.000794\pm0.001643$ and $-0.000496\pm0.002836$: no resolved final-checkpoint
advantage. After the 600-step tail the interaction comparison is modestly
favorable, $-0.002268\pm0.002154$, while the additive interval spans zero.
This isolated tail effect does not establish a general efficiency result.
Full-batch training on 2,048 noisy labels can overfit both larger fixed
networks and grown ones. The deployed checkpoint and ungated tail therefore
must not be conflated.

At a matched training/probe update proxy, an offline width-ten learner given
the full fitting-plus-audit label pool beats the residual controller on all
four tasks. On additive and interaction data, residual minus this comparator
is $+0.007953\pm0.002875$ and $+0.010387\pm0.002514$ respectively. The
corresponding fixed-pool risks are 0.13152 and 0.14028. These models use a
different minibatch schedule and receive future audit labels at the outset;
the result is an alternative end-of-lifetime resource allocation, not a
causal estimate of the effect of labels alone. It prevents claiming a
predictive-risk advantage at this total label and training budget. The fixed
comparator deploys ten units, however, whereas the adaptive models average
3.25 and 5.625; it does not dominate them simultaneously in inference capacity.

An additional, explicitly post-hoc diagnostic addresses this capacity mismatch.
For the 64 larger-audit residual runs, a fixed learner uses the same final
width and initialization as the existing matched-width comparator, but trains
from the whole label pool with the already specified minibatch schedule and
adaptive training/probe budget. It actually consumes all 133,120 distinct
labels in every run. Residual minus this same-width pool comparator is
$+0.003111\pm0.000767$ on null data, $+0.001266\pm0.000415$ on linear data,
$+0.003491\pm0.000860$ on additive data, and
$+0.002246\pm0.001887$ on interaction data. This removes the deployed-width
mismatch and reinforces the resource-allocation limitation. The comparator
still receives labels offline, uses a different optimizer schedule, and
knows final width retrospectively. It is a diagnostic added after review,
not a preregistered primary finding or an online dominance theorem.

\begin{figure}[tb]
\centering\includegraphics[width=\linewidth]{resource_trajectories.pdf}
\caption{Measured risk versus cumulative training and candidate-probe
parameter--example updates at the larger audit budget. Points summarize
16 trajectories. The fixed width-ten star uses the full label pool at the
residual controller's terminal update proxy. The horizontal axis omits audit
prediction work, matrix arithmetic, optimizer constants, and memory; these
are separately recorded and this plot is not a matched-FLOP comparison.}
\label{fig:resources}
\end{figure}

\paragraph{Rejection is partly a property of the chosen certificate.}
For the primary residual method, 253 small-audit proposals improve on both
incumbent and continuation by more than $\tau$ according to the independent
input evaluator; 237 are declined. At the larger audit size, 168 proposals
meet this evaluator criterion and 90 are declined. These counts are
diagnostic estimates, not exact population facts. The final histories differ,
so the two sets are not the same fixed proposal pool.

The smaller-audit additive task contains 125 evaluator-positive proposals,
all declined. Of these, 124 also have positive empirical cost-adjusted gains
against both comparators; 116 would fail even a bound consisting only of
the empirical-Bernstein linear term. The theoretical linear floor is
$14\log(960)/(3(n_A-1))$, approximately 0.007825 at $n_A=4096$ and
0.001956 at $n_A=16384$, before accounting for variance and the growth margin.
Consequently these failures cannot be presented as empirical confirmation
of a minimax information lower bound. They demonstrate conservatism of this
valid sufficient bound in the realized training regime.

Prediction movement and evaluator gain are separately logged for every
pair. On positive nonlinear proposals their ratios are often moderately
above the Bayes-refit identity $S=\Delta$, rather than uniformly exhibiting
extreme cancellation. Figure~\ref{fig:evidence} displays the full spread
including rejected positive proposals. A small gain alone does not diagnose
whether the proposal is weak, the empirical evidence is noisy, or the
implemented bound is unnecessarily loose.

\begin{figure}[tb]
\centering\includegraphics[width=\linewidth]{evidence_diagnostics.pdf}
\caption{Measured evidence diagnostics. Left: the finite-bound radius versus
evaluator gain over continuation for nonlinear residual proposals; the
mean-gate radius is diagnostic only. Middle: audit prediction movement versus
positive evaluator gain at the larger finite-bound audit; the dashed line is
the Bayes-refit identity, not a fitted neural law. Right: final risk under
small-audit finite bounds (blue), large-audit finite bounds (green), and the
large-audit empirical-mean heuristic (red). All input-population quantities
are Monte Carlo estimates; decisions use sampled audit labels.}
\label{fig:evidence}
\end{figure}

The heuristic mean gate accepts 97 residual growth events and reaches
mean widths 4.0625 and 6 on additive and interaction data, with risks
0.13629 and 0.14713. It provides an attainable but uncertified reference.
Twenty of its retained steps have a positive evaluator risk change, all
smaller than 0.001. No residual finite-bound retained step has an evaluator
risk increase. Neither observation replaces the population guarantee or
constitutes an exact test of its coverage; the evaluator still samples inputs.

\paragraph{The supplied-module evidence experiment.}
Separate from hidden-parameter training, 16,000 exact-distribution Monte
Carlo repetitions implement the tests in Theorems~\ref{thm:reset-cost} and
\ref{thm:shared-evidence}, with $K\in\{2,4,8,16\}$, $\varepsilon=0.2$,
and family error target 0.05. Signal, mixed, no-signal, and risk-neutral
boundary laws are retained. Negative-binomial waiting counts, multinomial
input counts, and binomial label counts are sufficient-statistic simulations
of ordinary iid sampling; they do not supply exact label expectations to
the tests. Module knots and amplitudes are nevertheless given, as stipulated
by the theorem, and no hidden parameters are trained.

At $K=16$ on the signal law, the reset procedure uses a mean 295,409 labels
versus the shared procedure's fixed 41,354, a factor 7.14 with these
conservative constants. Family correctness is 0.998 versus 1.000 over 1,000
repetitions; zero observed errors do not imply zero error probability.
At $K=2$ the shared sufficient bound costs more, 3,506 versus mean 2,952,
which is retained in the plot. Direct ReLU, Brier-risk, variance, and KL
identity checks agree to $5.6\cdot10^{-17}$. These validate the finite
realization and simulation; the lower-bound proof supplies the scaling claim.

\begin{figure}[tb]
\centering\includegraphics[width=\linewidth]{evidence_diagnostic.pdf}
\caption{Supplied-module audit experiment, not ordinary neural training.
Theoretical expected costs and lower bounds are distinguished from 1,000-run
Monte Carlo means. Error intervals retain uncertainty at zero observed
failures. The shared batch is valid because every compared local action is
precommitted, whereas the reset interface withholds off-action observations.}
\label{fig:module-evidence}
\end{figure}
